\itemhead{第3章 小结与复习 · 第4题（教材p.89）}{XJ-B1-C03-SU-P089-Y1284}
\label{XJ-B1-C03-SU-P089-Y1284}
{\footnotesize 湘教版教材p.89 / PDF第96页 / 第3章 小结与复习 / 第4题}
\field{原题}{如何判断一个函数的奇偶性？你能解释奇函数、偶函数图象的对称性吗？你能用数学语言来表述函数的奇偶性吗？}
\field{参考答案}{先查定义域对称性，再化简 $f(-x)$ 与 $f(x)$ 比较；偶函数图象关于 $y$ 轴对称，奇函数图象关于原点对称。只核对一对点不能证明恒等关系。}
{\small \field{关联知识}{F08 奇偶性、对称与函数性质应用（教材帮p.113、114、115 / PDF122、123、124）}
\field{方法与易错点}{检查定义域，再代入 $-x$；求另一侧解析式时同时写出 $-x$ 所在区间。 易错：只看公式不看定义域；用一对点证明奇偶；负半轴代换遗漏最外层负号。}
\field{答案依据}{答案册未提供本栏目的独立对应答案；依据原题、教材概念补写参考回应。开放性总结不设唯一数值答案。}}
